প্রথম-ক্রমের ভাষার \emph{অর্থতত্ত্ব} সেই ভাষার একটি \emph{গঠন}
দিয়ে নির্ধারিত হয়। গঠনটিতে একটি \emph{সংজ্ঞাক্ষেত্র} থাকে এবং
প্রতিটি ধ্রুবক সংকেতের জন্য সেই ক্ষেত্রের একটি করে উপাদান বরাদ্দ
হয়। অপেক্ষক সংকেতের ব্যাখ্যা হয় সংজ্ঞাক্ষেত্রের উপর অপেক্ষক
দিয়ে, আর সম্পর্ক সংকেতের ব্যাখ্যা হয় সম্পর্ক দিয়ে। চলরাশির
সেট থেকে সংজ্ঞাক্ষেত্রে একটি অপেক্ষক হলো \emph{চলরাশি-আরোপ}।
\emph{পরিতৃপ্তি} সম্পর্ক গঠন, চলরাশি-আরোপ ও সূত্রকে যুক্ত করে।
$\Sat{M}{!A}[s]$ সংজ্ঞায়িত হয়~$!A$-এর গঠনের উপর আরোহ করে।
$\Sat{M}{!A}[s]$ কেবল~$!A$-তে বাস্তবে সংঘটিত সংকেতগুলির
ব্যাখ্যার উপর নির্ভর করে। বিশেষত, $!A$-তে মুক্ত চলরাশি না থাকলে
এটি $s$-এর উপর নির্ভর করে না। তাই $!A$ বাক্য হলে
$\Sat{M}{!A}$ হয়, যদি $\Sat{M}{!A}[s]$ হয় যেকোনো (বা প্রত্যেক)~$s$-এর জন্য।

পরিতৃপ্তি সম্পর্কই সব অর্থগত ধারণার ভিত্তি। কোনো বাক্য প্রত্যেক
গঠনে পরিতৃপ্ত হলে সেটি \emph{বৈধ}, $\Sat{{}}{!A}$। কোনো
বাক্যসমষ্টি~$\Gamma$ থেকে বাক্য $!A$ \emph{অনুসৃত হয়},
$\Gamma \Entails !A$, ঠিক তখনই, যখন~$\Gamma$-র প্রত্যেক বাক্যকে
পরিতৃপ্ত করে এমন সব~$\Struct{M}$-এর জন্য $\Sat{M}{!A}$ হয়।
এমন কোনো গঠন থাকলে যাতে সেট~$\Gamma$-র প্রত্যেক বাক্য পরিতৃপ্ত
হয়,~$\Gamma$ \emph{পরিতৃপ্তিযোগ্য}, অন্যথায়
\emph{অপরিতৃপ্তিযোগ্য}। এই ধারণাগুলি পরস্পর সম্পর্কিত। যেমন,
$\Gamma \Entails !A$ হয় ঠিক তখনই, যখন $\Gamma \cup \{\lnot
!A\}$ অপরিতৃপ্তিযোগ্য।
